\documentclass[10pt,a4paper]{amsart}
\usepackage[margin=19mm]{geometry}
\usepackage{amsmath,amssymb,mathtools}
\usepackage[scr=rsfs,cal=euler]{mathalfa}
\usepackage{xfrac}
\usepackage[unicode,hidelinks,bookmarks=false]{hyperref}
\newcommand{\ie}{i.e.\ }
\newcommand{\cf}{cf.\ }
\newcommand{\resp}{respectively\ }
\newcommand{\Subsection}{\subsection}
\providecommand{\nobreakdash}{\nobreak\hbox{-}}
\setlength{\emergencystretch}{2em}
\multlinegap0pt
\usepackage{bm}
\usepackage{bbm}
\usepackage{dsfont}
\usepackage{xspace}
\usepackage{cancel}
\usepackage{mathtools}
\usepackage{amsthm}
\makeatletter
\@ifundefined{thm}{\newtheorem{thm}{Thm}}{}
\@ifundefined{Proposition}{\newtheorem{Proposition}[thm]{Proposition}}{}
\makeatother
\title[Linearizations and optimization problems in diffeological sp]{Linearizations and optimization problems in diffeological spaces}
\author{Jean-Pierre Magnot}
\date{Source: arXiv:2507.19508v2 (CC BY 4.0)}
\begin{document}
\maketitle

\noindent\emph{Source and licence.} Linearizations and optimization problems in diffeological spaces. Version arXiv:2507.19508v2 (CC BY 4.0), distributed under the Creative Commons Attribution 4.0 International licence, as recorded in the arXiv licence metadata for this exact version and shown on the abstract page https://arxiv.org/abs/2507.19508v2. Full rights notice: licenses/arxiv-2507.19508-CC-BY-4.0.txt.\par\smallskip

\noindent\textit{Editorial scope.} Counted unit: the Proposition whose source label is \texttt{prop:linearization-smooth} in the article (source0.tex lines 409--442, source label \texttt{prop:linearization-smooth}), delivered together with the article's own complete original proof of it (source0.tex lines 444--588, one \texttt{proof} environment of 145 lines). No mathematics is written, repaired or re-derived: statement and proof are the article's own text line for line. Required context retained uncounted: none. Every definition, display and label of those lines is retained unchanged. Omitted from this package and not redistributed: 1--408, 443--443, 589--1480. Nothing inside a retained excerpt is omitted or altered: every retained body is delivered exactly as the article's lines are written, so no omission receipt of any kind accompanies this package. Citations: the retained excerpts carry no citation command at all, so no bibliography entry of the article is transcribed and no reference is redistributed. Reference adaptation: none was needed and none was made; every reference inside the delivered unit resolves inside it, so no unresolved reference or citation is produced. Printed slips kept literal: no printed word, symbol, hypothesis or number of the counted statement or of its proof was altered. Layout and class: the article's own class is \texttt{amsart} with packages including inputenc, amsmath, amssymb, amsthm, babel, xcolor, pgf, tikz-cd, algorithmic, algorithm2e; the wrapper is the lane's pinned \texttt{amsart} editorial wrapper at 10pt on A4, which restates the theorem-like environments the retained text needs and adds the article's own macro definitions that the retained text needs (none), with extra wrapper packages (bm, bbm, dsfont, xspace, cancel, mathtools). The delivered numbering is the unit's own. Assets: no retained span contains a figure environment, a graphics-inclusion command, a PDF-page inclusion, a table or a TikZ picture; no image file is copied into this package, the package declares no asset dependency and no image file is redistributed with it. Licence: version arXiv:2507.19508v2 of this article is distributed under the Creative Commons Attribution 4.0 International licence, as recorded in the arXiv licence metadata for this exact version and shown on the abstract page https://arxiv.org/abs/2507.19508v2. A full rights notice is delivered at licenses/arxiv-2507.19508-CC-BY-4.0.txt. Characters: every retained excerpt is pure ASCII, so the delivered engine, XeLaTeX with its default Latin Modern text font, reports no missing character.
\begin{Proposition}[Smooth linearization on a differentiable manifold]
\label{prop:linearization-smooth}
Let $(X,g)$ be a smooth finite-dimensional Riemannian manifold.
Using the musical isomorphism $v \mapsto v^\flat := g_x(v,\cdot)$,
we identify $T_xX$ with $T_x^\ast X$.
Set
\[
E := T^\ast X \oplus \mathbb{R},
\]
viewed as a smooth vector bundle over $X$ via the projection 
$\pi_1 : E \to X$.

Then there exists a smooth map
\[
\nu : X\times X \longrightarrow E
\]
such that:
\begin{enumerate}
  \item[(i)] $\nu(x,x) = (0_x,0)$ for all $x\in X$;

  \item[(ii)] $\nu(x,y) = (0_x,0)$ if and only if $y=x$;

  \item[(iii)] for each $x\in X$, the differential of the first component
  of $\nu$ with respect to the second variable, evaluated at $(x,x)$, is
  the musical isomorphism
  \[
    D_2\bigl(pr_1\circ\nu\bigr)(x,x)
    = (\cdot)^\flat : T_xX \longrightarrow T_x^\ast X ,
  \]
  where $pr_1 : E\to T^\ast X$ denotes the first projection.
\end{enumerate}
In particular, $(E,\nu)$ defines a smooth linearization of $X$ along the
diagonal in the above sense.
\end{Proposition}

\begin{proof}
\emph{Step 1: A Bokobza--Haggiag type linear part in $T^\ast X$.}

Fix a smooth Riemannian metric $g$ on $X$.  
For each $x\in X$, let $\exp_x : \mathcal U_x \subset T_xX \to X$ be the 
exponential map, defined on some neighbourhood $\mathcal U_x$ of $0_x$.
By the tubular neighbourhood theorem for the diagonal
\[
\Delta := \{(x,x) : x\in X\} \subset X\times X,
\]
there exists an open neighbourhood $U$ of $\Delta$ in $X\times X$ and
a smooth map
\[
\xi : U \longrightarrow TX
\]
such that:
\begin{enumerate}
  \item $\xi(x,y)\in T_xX$ for all $(x,y)\in U$;
  \item $\xi(x,x) = 0_x$ for all $x\in X$;
  \item $y = \exp_x\bigl(\xi(x,y)\bigr)$ for all $(x,y)\in U$.
\end{enumerate}
In particular,
\[
D_2\xi(x,x) = \mathrm{Id}_{T_xX}.
\]

Define
\[
\nu_{\mathrm{BH}} : U \longrightarrow T^\ast X,
\qquad
\nu_{\mathrm{BH}}(x,y) := \xi(x,y)^\flat.
\]
Then $\nu_{\mathrm{BH}}$ is smooth on $U$, satisfies
$\nu_{\mathrm{BH}}(x,x)=0_x$, and
\[
D_2\nu_{\mathrm{BH}}(x,x)
= (\cdot)^\flat : T_xX \to T_x^\ast X.
\]

\medskip
\noindent
\emph{Step 2: A smooth function vanishing exactly on the diagonal.}

Define on $U$:
\[
\rho_0(x,y) := \|\xi(x,y)\|_g^2.
\]
Then $\rho_0$ is smooth on $U$, nonnegative, and
\[
\rho_0(x,y) = 0 \Longleftrightarrow (x,y)\in \Delta \cap U.
\]

Choose an open subset $U_0$ such that
\[
\Delta \subset U_0 \subset \overline{U_0} \subset U,
\]
and a smooth cutoff function 
$\chi\in C^\infty(X\times X,[0,1])$ with
\[
\chi \equiv 1 \ \text{on } U_0,
\qquad
\mathrm{supp}(\chi)\subset U.
\]
Define a global smooth map
\[
\rho : X\times X \longrightarrow [0,+\infty)
\]
by
\[
\rho(x,y)
:= \chi(x,y)\,\rho_0(x,y) + \bigl(1-\chi(x,y)\bigr).
\]
Then:
\begin{itemize}
  \item on $U_0$, $\rho=\rho_0$, hence $\rho(x,y)=0$ iff $(x,y)\in\Delta$;
  \item outside $U$, $\chi=0$ and thus $\rho\equiv 1$.
\end{itemize}
Hence $\rho$ is smooth and
\[
\rho(x,y)=0 \Longleftrightarrow (x,y)\in\Delta.
\]

\medskip
\noindent
\emph{Step 3: Global extension of $\nu_{\mathrm{BH}}$.}

Define
\[
\widetilde{\nu}_{\mathrm{BH}}(x,y)
:= \chi(x,y)\,\nu_{\mathrm{BH}}(x,y),
\]
extending $\nu_{\mathrm{BH}}$ smoothly across $U$ and taking it to $0$ outside $U$.
Thus $\widetilde{\nu}_{\mathrm{BH}}$ is smooth on $X\times X$, and coincides
with $\nu_{\mathrm{BH}}$ on $U_0$.

\medskip
\noindent
\emph{Step 4: Definition of the linearization $\nu$.}

Set
\[
\nu(x,y)
:= \bigl(\widetilde{\nu}_{\mathrm{BH}}(x,y),\;\rho(x,y)\bigr)
\in T_x^\ast X \oplus \mathbb{R}.
\]

\medskip
\noindent
\emph{Step 5: Verification of properties (i)--(iii).}

(i)  
For $y=x$, we have $(x,x)\in U_0$.
Hence $\chi(x,x)=1$, $\rho(x,x)=\|\xi(x,x)\|^2=0$, and
$\widetilde{\nu}_{\mathrm{BH}}(x,x)=\nu_{\mathrm{BH}}(x,x)=0_x$.
Thus
\[
\nu(x,x)=(0_x,0).
\]

(ii)  
If $\nu(x,y)=(0_x,0)$, then in particular $\rho(x,y)=0$.
By construction of $\rho$, this implies $(x,y)\in\Delta$, hence $y=x$.
The converse is already proved in (i).
Thus
\[
\nu(x,y)=(0_x,0)\quad\Longleftrightarrow\quad y=x.
\]

(iii)  
Near $(x,x)$, we are in $U_0$, so
\[
pr_1\circ \nu = \widetilde{\nu}_{\mathrm{BH}}
                = \nu_{\mathrm{BH}}.
\]
Therefore
\[
D_2(pr_1\circ \nu)(x,x)
= D_2\nu_{\mathrm{BH}}(x,x)
= (\cdot)^\flat : T_xX \to T_x^\ast X,
\]
which is the required linear isomorphism.

\medskip
This completes the proof.
\end{proof}

\end{document}
